%! Exponential Growth and Decay — Unit 7, Lesson 7.1 — Warm-Up
\documentclass[10pt]{article}
\usepackage{algebra2-article}
\usepackage{algebra2-boxes}
\newcommand{\TallMath}[1]{$\displaystyle #1\rule[-1.4em]{0pt}{3.2em}$}

\begin{document}

\pageheader{Unit 7, Lesson 7.1}{Warm-Up}
\namedateperiod

\vspace{0.10in}

\begin{notesbox}{Getting Ready: Skills We'll Use Today}
\small These three skills power today's lesson on \emph{building} an exponential model. Work quickly.

\vspace{4pt}
\textbf{1. Percent change --- the long way, then the short way.} A jacket costs \$$50$.

\vspace{3pt}
(a) The price goes \textbf{up} $8\%$. \\[2pt]
\hspace*{1.2em} The increase is \blank{1.6cm}, so the new price is \blank{1.8cm}. \\[2pt]
\hspace*{1.2em} Now get the same answer with \emph{one} multiplication:
$\$50 \times$ \blank{1.4cm} $=$ \blank{1.8cm}. \\[2pt]
\hspace*{1.2em} The new price is \blank{1.4cm}\% of the old price.

\vspace{4pt}
(b) Starting over from \$$50$, the price goes \textbf{down} $15\%$. \\[2pt]
\hspace*{1.2em} The decrease is \blank{1.6cm}, so the new price is \blank{1.8cm}. \\[2pt]
\hspace*{1.2em} One multiplication: $\$50 \times$ \blank{1.4cm} $=$ \blank{1.8cm}. \\[2pt]
\hspace*{1.2em} The new price is \blank{1.4cm}\% of the old price.

\vspace{5pt}
\textbf{2. Reading $a$ and $b$ (Lesson 7.0).} Fill in the table from the equation alone.
\begin{center}
\renewcommand{\arraystretch}{1.5}
\begin{tabularx}{\linewidth}{@{}l X X X X@{}}
\hline
\textbf{Function} & \textbf{$a$} & \textbf{$b$} & \textbf{Growth or decay?} & \textbf{$y$-intercept} \\
\hline
$y=5(1.2)^{x}$ & \blank{1.6cm} & \blank{1.6cm} & \blank{2.0cm} & \blank{1.8cm} \\
$y=80(0.6)^{x}$ & \blank{1.6cm} & \blank{1.6cm} & \blank{2.0cm} & \blank{1.8cm} \\
\hline
\end{tabularx}
\end{center}

\vspace{2pt}
\textbf{3. A constant ratio.} The table below is exponential.
\begin{center}
\renewcommand{\arraystretch}{1.3}
\begin{tabular}{|c|c|c|c|c|}
\hline
$x$ & $0$ & $1$ & $2$ & $3$ \\ \hline
$y$ & $7$ & $21$ & $63$ & $189$ \\ \hline
\end{tabular}
\end{center}
Divide to find the constant ratio: $\dfrac{21}{7}=$ \blank{1.0cm}, \; $\dfrac{63}{21}=$
\blank{1.0cm}, \; $\dfrac{189}{63}=$ \blank{1.0cm}. \\[3pt]
So the base is $b=$ \blank{1.2cm}, and the output at $x=0$ is $a=$ \blank{1.2cm}.

\end{notesbox}

\end{document}
